% =====================================================================
%  Standard-Präambel für Skripte und Aufgabensammlungen
%  Gymnasium Baden-Württemberg, Bildungsplan 2016
%
%  Verwendung:  % =====================================================================
%  Standard-Präambel für Skripte und Aufgabensammlungen
%  Gymnasium Baden-Württemberg, Bildungsplan 2016
%
%  Verwendung:  % =====================================================================
%  Standard-Präambel für Skripte und Aufgabensammlungen
%  Gymnasium Baden-Württemberg, Bildungsplan 2016
%
%  Verwendung:  % =====================================================================
%  Standard-Präambel für Skripte und Aufgabensammlungen
%  Gymnasium Baden-Württemberg, Bildungsplan 2016
%
%  Verwendung:  \input{latex-praeambel.tex}  oder Inhalt direkt einfügen
%  Engine:      pdflatex (zwei Durchläufe für Inhaltsverzeichnis)
%               XeLaTeX nur bei Bedarf (fontspec/Spezialfonts)
% =====================================================================

\documentclass[11pt,a4paper,parskip=half]{scrartcl}

% --- Sprache & Eingabe ----------------------------------------------
\usepackage[utf8]{inputenc}
\usepackage[T1]{fontenc}
\usepackage[ngerman]{babel}
\usepackage{lmodern}
\usepackage{microtype}

% --- Mathematik -----------------------------------------------------
\usepackage{amsmath,amssymb,amsthm,mathtools}
\usepackage{cancel}                  % \cancel für Kürzen
\usepackage{siunitx}                 % \SI, \num, \si — deutsche Konvention
\sisetup{
  output-decimal-marker = {,},       % Komma statt Punkt
  per-mode              = symbol,    % m/s statt m·s^{-1}
  range-phrase          = {\,bis\,},
  list-final-separator  = { und },
  group-separator       = {\,},      % 10 000 mit schmalem Leerzeichen
}

% --- Layout ---------------------------------------------------------
\usepackage[a4paper,left=2.4cm,right=2.4cm,top=2.5cm,bottom=2.5cm]{geometry}
\usepackage{enumitem}
\setlist{topsep=2pt,itemsep=2pt}
\usepackage{array,booktabs,longtable,tabularx,multicol}
\usepackage{ragged2e}
\usepackage{fancyhdr}
\usepackage{hyperref}
\hypersetup{
  colorlinks = true,
  linkcolor  = mainblue,
  urlcolor   = mainblue,
  citecolor  = mainblue,
}

% --- Grafik ---------------------------------------------------------
\usepackage{tikz}
\usepackage{pgfplots}
\pgfplotsset{compat=1.18}
\usetikzlibrary{
  patterns, arrows.meta, calc, positioning,
  decorations.pathreplacing, decorations.markings,
  intersections, angles, quotes, 3d
}

% --- Chemie (nur einbinden, wenn benötigt) --------------------------
% \usepackage[version=4]{mhchem}     % Reaktionsgleichungen: \ce{H2O}
% \usepackage{chemfig}               % Strukturformeln
% \setatomsep{1.8em}
% \setbondoffset{1pt}
% \setdoublesep{0.45ex}

% =====================================================================
%  Farbschema (Kleeblatt-Konvention für Skripte)
% =====================================================================
\usepackage{xcolor}
\definecolor{mainblue}   {RGB}{ 0, 90,160}   % Definitionen, Hauptfarbe
\definecolor{darkblue}   {RGB}{ 0, 50,100}
\definecolor{exgreen}    {RGB}{ 34,120, 67}  % Beispiele
\definecolor{aufgorange} {RGB}{210,120, 20}  % Aufgaben
\definecolor{loesgray}   {RGB}{ 90, 90, 90}  % Lösungen
\definecolor{hinwyellow} {RGB}{200,160,  0}  % Hinweise
\definecolor{stolperred} {RGB}{180, 40, 40}  % Stolperfallen
\definecolor{tipppurple} {RGB}{120, 40,140}  % Klausurtipps
\definecolor{lighgray}   {RGB}{245,245,245}

% =====================================================================
%  Boxen-Umgebungen (tcolorbox)
% =====================================================================
\usepackage[most]{tcolorbox}
\tcbuselibrary{theorems, skins, breakable}

% --- Definition / Merksatz ------------------------------------------
\newtcolorbox{definition}[1][Definition]{%
  enhanced, breakable,
  colback=mainblue!7, colframe=mainblue,
  fonttitle=\bfseries,
  title=#1,
  attach boxed title to top left={yshift=-2mm,xshift=4mm},
  boxed title style={colback=mainblue,colframe=mainblue},
  arc=2pt, top=3mm, bottom=2mm
}

% --- Beispiel mit Musterlösung --------------------------------------
\newtcolorbox{beispiel}[1][Beispiel]{%
  enhanced, breakable,
  colback=exgreen!7, colframe=exgreen,
  fonttitle=\bfseries,
  title=#1,
  attach boxed title to top left={yshift=-2mm,xshift=4mm},
  boxed title style={colback=exgreen,colframe=exgreen},
  arc=2pt, top=3mm, bottom=2mm
}

% --- Übungsaufgabe --------------------------------------------------
\newtcolorbox{aufgabe}[1][Aufgabe]{%
  enhanced, breakable,
  colback=aufgorange!8, colframe=aufgorange,
  fonttitle=\bfseries,
  title=#1,
  attach boxed title to top left={yshift=-2mm,xshift=4mm},
  boxed title style={colback=aufgorange,colframe=aufgorange},
  arc=2pt, top=3mm, bottom=2mm
}

% --- Lösung ---------------------------------------------------------
\newtcolorbox{loesung}[1][Lösung]{%
  enhanced, breakable,
  colback=lighgray, colframe=loesgray,
  fonttitle=\bfseries,
  title=#1,
  attach boxed title to top left={yshift=-2mm,xshift=4mm},
  boxed title style={colback=loesgray,colframe=loesgray},
  arc=2pt, top=3mm, bottom=2mm
}

% --- Allgemeiner Hinweis --------------------------------------------
\newtcolorbox{hinweis}[1][Hinweis]{%
  enhanced, breakable,
  colback=hinwyellow!10, colframe=hinwyellow,
  fonttitle=\bfseries,
  title=#1,
  attach boxed title to top left={yshift=-2mm,xshift=4mm},
  boxed title style={colback=hinwyellow,colframe=hinwyellow},
  arc=2pt, top=3mm, bottom=2mm
}

% --- Stolperfalle (typischer Klausurfehler) -------------------------
\newtcolorbox{stolperfalle}[1][Stolperfalle]{%
  enhanced, breakable,
  colback=stolperred!8, colframe=stolperred,
  fonttitle=\bfseries\color{white},
  title=#1,
  attach boxed title to top left={yshift=-2mm,xshift=4mm},
  boxed title style={colback=stolperred,colframe=stolperred},
  arc=2pt, top=3mm, bottom=2mm
}

% --- Klausurtipp ----------------------------------------------------
\newtcolorbox{klausurtipp}[1][Klausurtipp]{%
  enhanced, breakable,
  colback=tipppurple!8, colframe=tipppurple,
  fonttitle=\bfseries\color{white},
  title=#1,
  attach boxed title to top left={yshift=-2mm,xshift=4mm},
  boxed title style={colback=tipppurple,colframe=tipppurple},
  arc=2pt, top=3mm, bottom=2mm
}

% =====================================================================
%  Sonstige Konventionen
% =====================================================================
\renewcommand{\vec}[1]{\boldsymbol{#1}}        % Vektoren fett statt Pfeil
% Alternative: \renewcommand{\vec}[1]{\overrightarrow{#1}}

\newcommand{\R}{\mathbb{R}}
\newcommand{\N}{\mathbb{N}}
\newcommand{\Z}{\mathbb{Z}}
\newcommand{\Q}{\mathbb{Q}}

% Endergebnis einrahmen (für Aufgabenlösungen)
\newcommand{\ergebnis}[1]{\boxed{\;#1\;}}

% Kurzbefehl für Bildungsplan-Code-Verweise
\newcommand{\bpref}[1]{\textit{(Bildungsplan #1)}}

% =====================================================================
%  Kopf- und Fußzeile (optional, hier aktiv)
% =====================================================================
\pagestyle{fancy}
\fancyhf{}
\fancyhead[L]{\small\itshape\leftmark}
\fancyhead[R]{\small Klasse 10 -- Gymnasium BW}
\fancyfoot[C]{\thepage}
\renewcommand{\headrulewidth}{0.4pt}

% =====================================================================
%  Beginn des Dokuments
% =====================================================================
\begin{document}

% \title{...}
% \author{...}
% \date{\today}
% \maketitle
% \tableofcontents
% \newpage

% --- Inhalt hier einfügen ---

\end{document}
  oder Inhalt direkt einfügen
%  Engine:      pdflatex (zwei Durchläufe für Inhaltsverzeichnis)
%               XeLaTeX nur bei Bedarf (fontspec/Spezialfonts)
% =====================================================================

\documentclass[11pt,a4paper,parskip=half]{scrartcl}

% --- Sprache & Eingabe ----------------------------------------------
\usepackage[utf8]{inputenc}
\usepackage[T1]{fontenc}
\usepackage[ngerman]{babel}
\usepackage{lmodern}
\usepackage{microtype}

% --- Mathematik -----------------------------------------------------
\usepackage{amsmath,amssymb,amsthm,mathtools}
\usepackage{cancel}                  % \cancel für Kürzen
\usepackage{siunitx}                 % \SI, \num, \si — deutsche Konvention
\sisetup{
  output-decimal-marker = {,},       % Komma statt Punkt
  per-mode              = symbol,    % m/s statt m·s^{-1}
  range-phrase          = {\,bis\,},
  list-final-separator  = { und },
  group-separator       = {\,},      % 10 000 mit schmalem Leerzeichen
}

% --- Layout ---------------------------------------------------------
\usepackage[a4paper,left=2.4cm,right=2.4cm,top=2.5cm,bottom=2.5cm]{geometry}
\usepackage{enumitem}
\setlist{topsep=2pt,itemsep=2pt}
\usepackage{array,booktabs,longtable,tabularx,multicol}
\usepackage{ragged2e}
\usepackage{fancyhdr}
\usepackage{hyperref}
\hypersetup{
  colorlinks = true,
  linkcolor  = mainblue,
  urlcolor   = mainblue,
  citecolor  = mainblue,
}

% --- Grafik ---------------------------------------------------------
\usepackage{tikz}
\usepackage{pgfplots}
\pgfplotsset{compat=1.18}
\usetikzlibrary{
  patterns, arrows.meta, calc, positioning,
  decorations.pathreplacing, decorations.markings,
  intersections, angles, quotes, 3d
}

% --- Chemie (nur einbinden, wenn benötigt) --------------------------
% \usepackage[version=4]{mhchem}     % Reaktionsgleichungen: \ce{H2O}
% \usepackage{chemfig}               % Strukturformeln
% \setatomsep{1.8em}
% \setbondoffset{1pt}
% \setdoublesep{0.45ex}

% =====================================================================
%  Farbschema (Kleeblatt-Konvention für Skripte)
% =====================================================================
\usepackage{xcolor}
\definecolor{mainblue}   {RGB}{ 0, 90,160}   % Definitionen, Hauptfarbe
\definecolor{darkblue}   {RGB}{ 0, 50,100}
\definecolor{exgreen}    {RGB}{ 34,120, 67}  % Beispiele
\definecolor{aufgorange} {RGB}{210,120, 20}  % Aufgaben
\definecolor{loesgray}   {RGB}{ 90, 90, 90}  % Lösungen
\definecolor{hinwyellow} {RGB}{200,160,  0}  % Hinweise
\definecolor{stolperred} {RGB}{180, 40, 40}  % Stolperfallen
\definecolor{tipppurple} {RGB}{120, 40,140}  % Klausurtipps
\definecolor{lighgray}   {RGB}{245,245,245}

% =====================================================================
%  Boxen-Umgebungen (tcolorbox)
% =====================================================================
\usepackage[most]{tcolorbox}
\tcbuselibrary{theorems, skins, breakable}

% --- Definition / Merksatz ------------------------------------------
\newtcolorbox{definition}[1][Definition]{%
  enhanced, breakable,
  colback=mainblue!7, colframe=mainblue,
  fonttitle=\bfseries,
  title=#1,
  attach boxed title to top left={yshift=-2mm,xshift=4mm},
  boxed title style={colback=mainblue,colframe=mainblue},
  arc=2pt, top=3mm, bottom=2mm
}

% --- Beispiel mit Musterlösung --------------------------------------
\newtcolorbox{beispiel}[1][Beispiel]{%
  enhanced, breakable,
  colback=exgreen!7, colframe=exgreen,
  fonttitle=\bfseries,
  title=#1,
  attach boxed title to top left={yshift=-2mm,xshift=4mm},
  boxed title style={colback=exgreen,colframe=exgreen},
  arc=2pt, top=3mm, bottom=2mm
}

% --- Übungsaufgabe --------------------------------------------------
\newtcolorbox{aufgabe}[1][Aufgabe]{%
  enhanced, breakable,
  colback=aufgorange!8, colframe=aufgorange,
  fonttitle=\bfseries,
  title=#1,
  attach boxed title to top left={yshift=-2mm,xshift=4mm},
  boxed title style={colback=aufgorange,colframe=aufgorange},
  arc=2pt, top=3mm, bottom=2mm
}

% --- Lösung ---------------------------------------------------------
\newtcolorbox{loesung}[1][Lösung]{%
  enhanced, breakable,
  colback=lighgray, colframe=loesgray,
  fonttitle=\bfseries,
  title=#1,
  attach boxed title to top left={yshift=-2mm,xshift=4mm},
  boxed title style={colback=loesgray,colframe=loesgray},
  arc=2pt, top=3mm, bottom=2mm
}

% --- Allgemeiner Hinweis --------------------------------------------
\newtcolorbox{hinweis}[1][Hinweis]{%
  enhanced, breakable,
  colback=hinwyellow!10, colframe=hinwyellow,
  fonttitle=\bfseries,
  title=#1,
  attach boxed title to top left={yshift=-2mm,xshift=4mm},
  boxed title style={colback=hinwyellow,colframe=hinwyellow},
  arc=2pt, top=3mm, bottom=2mm
}

% --- Stolperfalle (typischer Klausurfehler) -------------------------
\newtcolorbox{stolperfalle}[1][Stolperfalle]{%
  enhanced, breakable,
  colback=stolperred!8, colframe=stolperred,
  fonttitle=\bfseries\color{white},
  title=#1,
  attach boxed title to top left={yshift=-2mm,xshift=4mm},
  boxed title style={colback=stolperred,colframe=stolperred},
  arc=2pt, top=3mm, bottom=2mm
}

% --- Klausurtipp ----------------------------------------------------
\newtcolorbox{klausurtipp}[1][Klausurtipp]{%
  enhanced, breakable,
  colback=tipppurple!8, colframe=tipppurple,
  fonttitle=\bfseries\color{white},
  title=#1,
  attach boxed title to top left={yshift=-2mm,xshift=4mm},
  boxed title style={colback=tipppurple,colframe=tipppurple},
  arc=2pt, top=3mm, bottom=2mm
}

% =====================================================================
%  Sonstige Konventionen
% =====================================================================
\renewcommand{\vec}[1]{\boldsymbol{#1}}        % Vektoren fett statt Pfeil
% Alternative: \renewcommand{\vec}[1]{\overrightarrow{#1}}

\newcommand{\R}{\mathbb{R}}
\newcommand{\N}{\mathbb{N}}
\newcommand{\Z}{\mathbb{Z}}
\newcommand{\Q}{\mathbb{Q}}

% Endergebnis einrahmen (für Aufgabenlösungen)
\newcommand{\ergebnis}[1]{\boxed{\;#1\;}}

% Kurzbefehl für Bildungsplan-Code-Verweise
\newcommand{\bpref}[1]{\textit{(Bildungsplan #1)}}

% =====================================================================
%  Kopf- und Fußzeile (optional, hier aktiv)
% =====================================================================
\pagestyle{fancy}
\fancyhf{}
\fancyhead[L]{\small\itshape\leftmark}
\fancyhead[R]{\small Klasse 10 -- Gymnasium BW}
\fancyfoot[C]{\thepage}
\renewcommand{\headrulewidth}{0.4pt}

% =====================================================================
%  Beginn des Dokuments
% =====================================================================
\begin{document}

% \title{...}
% \author{...}
% \date{\today}
% \maketitle
% \tableofcontents
% \newpage

% --- Inhalt hier einfügen ---

\end{document}
  oder Inhalt direkt einfügen
%  Engine:      pdflatex (zwei Durchläufe für Inhaltsverzeichnis)
%               XeLaTeX nur bei Bedarf (fontspec/Spezialfonts)
% =====================================================================

\documentclass[11pt,a4paper,parskip=half]{scrartcl}

% --- Sprache & Eingabe ----------------------------------------------
\usepackage[utf8]{inputenc}
\usepackage[T1]{fontenc}
\usepackage[ngerman]{babel}
\usepackage{lmodern}
\usepackage{microtype}

% --- Mathematik -----------------------------------------------------
\usepackage{amsmath,amssymb,amsthm,mathtools}
\usepackage{cancel}                  % \cancel für Kürzen
\usepackage{siunitx}                 % \SI, \num, \si — deutsche Konvention
\sisetup{
  output-decimal-marker = {,},       % Komma statt Punkt
  per-mode              = symbol,    % m/s statt m·s^{-1}
  range-phrase          = {\,bis\,},
  list-final-separator  = { und },
  group-separator       = {\,},      % 10 000 mit schmalem Leerzeichen
}

% --- Layout ---------------------------------------------------------
\usepackage[a4paper,left=2.4cm,right=2.4cm,top=2.5cm,bottom=2.5cm]{geometry}
\usepackage{enumitem}
\setlist{topsep=2pt,itemsep=2pt}
\usepackage{array,booktabs,longtable,tabularx,multicol}
\usepackage{ragged2e}
\usepackage{fancyhdr}
\usepackage{hyperref}
\hypersetup{
  colorlinks = true,
  linkcolor  = mainblue,
  urlcolor   = mainblue,
  citecolor  = mainblue,
}

% --- Grafik ---------------------------------------------------------
\usepackage{tikz}
\usepackage{pgfplots}
\pgfplotsset{compat=1.18}
\usetikzlibrary{
  patterns, arrows.meta, calc, positioning,
  decorations.pathreplacing, decorations.markings,
  intersections, angles, quotes, 3d
}

% --- Chemie (nur einbinden, wenn benötigt) --------------------------
% \usepackage[version=4]{mhchem}     % Reaktionsgleichungen: \ce{H2O}
% \usepackage{chemfig}               % Strukturformeln
% \setatomsep{1.8em}
% \setbondoffset{1pt}
% \setdoublesep{0.45ex}

% =====================================================================
%  Farbschema (Kleeblatt-Konvention für Skripte)
% =====================================================================
\usepackage{xcolor}
\definecolor{mainblue}   {RGB}{ 0, 90,160}   % Definitionen, Hauptfarbe
\definecolor{darkblue}   {RGB}{ 0, 50,100}
\definecolor{exgreen}    {RGB}{ 34,120, 67}  % Beispiele
\definecolor{aufgorange} {RGB}{210,120, 20}  % Aufgaben
\definecolor{loesgray}   {RGB}{ 90, 90, 90}  % Lösungen
\definecolor{hinwyellow} {RGB}{200,160,  0}  % Hinweise
\definecolor{stolperred} {RGB}{180, 40, 40}  % Stolperfallen
\definecolor{tipppurple} {RGB}{120, 40,140}  % Klausurtipps
\definecolor{lighgray}   {RGB}{245,245,245}

% =====================================================================
%  Boxen-Umgebungen (tcolorbox)
% =====================================================================
\usepackage[most]{tcolorbox}
\tcbuselibrary{theorems, skins, breakable}

% --- Definition / Merksatz ------------------------------------------
\newtcolorbox{definition}[1][Definition]{%
  enhanced, breakable,
  colback=mainblue!7, colframe=mainblue,
  fonttitle=\bfseries,
  title=#1,
  attach boxed title to top left={yshift=-2mm,xshift=4mm},
  boxed title style={colback=mainblue,colframe=mainblue},
  arc=2pt, top=3mm, bottom=2mm
}

% --- Beispiel mit Musterlösung --------------------------------------
\newtcolorbox{beispiel}[1][Beispiel]{%
  enhanced, breakable,
  colback=exgreen!7, colframe=exgreen,
  fonttitle=\bfseries,
  title=#1,
  attach boxed title to top left={yshift=-2mm,xshift=4mm},
  boxed title style={colback=exgreen,colframe=exgreen},
  arc=2pt, top=3mm, bottom=2mm
}

% --- Übungsaufgabe --------------------------------------------------
\newtcolorbox{aufgabe}[1][Aufgabe]{%
  enhanced, breakable,
  colback=aufgorange!8, colframe=aufgorange,
  fonttitle=\bfseries,
  title=#1,
  attach boxed title to top left={yshift=-2mm,xshift=4mm},
  boxed title style={colback=aufgorange,colframe=aufgorange},
  arc=2pt, top=3mm, bottom=2mm
}

% --- Lösung ---------------------------------------------------------
\newtcolorbox{loesung}[1][Lösung]{%
  enhanced, breakable,
  colback=lighgray, colframe=loesgray,
  fonttitle=\bfseries,
  title=#1,
  attach boxed title to top left={yshift=-2mm,xshift=4mm},
  boxed title style={colback=loesgray,colframe=loesgray},
  arc=2pt, top=3mm, bottom=2mm
}

% --- Allgemeiner Hinweis --------------------------------------------
\newtcolorbox{hinweis}[1][Hinweis]{%
  enhanced, breakable,
  colback=hinwyellow!10, colframe=hinwyellow,
  fonttitle=\bfseries,
  title=#1,
  attach boxed title to top left={yshift=-2mm,xshift=4mm},
  boxed title style={colback=hinwyellow,colframe=hinwyellow},
  arc=2pt, top=3mm, bottom=2mm
}

% --- Stolperfalle (typischer Klausurfehler) -------------------------
\newtcolorbox{stolperfalle}[1][Stolperfalle]{%
  enhanced, breakable,
  colback=stolperred!8, colframe=stolperred,
  fonttitle=\bfseries\color{white},
  title=#1,
  attach boxed title to top left={yshift=-2mm,xshift=4mm},
  boxed title style={colback=stolperred,colframe=stolperred},
  arc=2pt, top=3mm, bottom=2mm
}

% --- Klausurtipp ----------------------------------------------------
\newtcolorbox{klausurtipp}[1][Klausurtipp]{%
  enhanced, breakable,
  colback=tipppurple!8, colframe=tipppurple,
  fonttitle=\bfseries\color{white},
  title=#1,
  attach boxed title to top left={yshift=-2mm,xshift=4mm},
  boxed title style={colback=tipppurple,colframe=tipppurple},
  arc=2pt, top=3mm, bottom=2mm
}

% =====================================================================
%  Sonstige Konventionen
% =====================================================================
\renewcommand{\vec}[1]{\boldsymbol{#1}}        % Vektoren fett statt Pfeil
% Alternative: \renewcommand{\vec}[1]{\overrightarrow{#1}}

\newcommand{\R}{\mathbb{R}}
\newcommand{\N}{\mathbb{N}}
\newcommand{\Z}{\mathbb{Z}}
\newcommand{\Q}{\mathbb{Q}}

% Endergebnis einrahmen (für Aufgabenlösungen)
\newcommand{\ergebnis}[1]{\boxed{\;#1\;}}

% Kurzbefehl für Bildungsplan-Code-Verweise
\newcommand{\bpref}[1]{\textit{(Bildungsplan #1)}}

% =====================================================================
%  Kopf- und Fußzeile (optional, hier aktiv)
% =====================================================================
\pagestyle{fancy}
\fancyhf{}
\fancyhead[L]{\small\itshape\leftmark}
\fancyhead[R]{\small Klasse 10 -- Gymnasium BW}
\fancyfoot[C]{\thepage}
\renewcommand{\headrulewidth}{0.4pt}

% =====================================================================
%  Beginn des Dokuments
% =====================================================================
\begin{document}

% \title{...}
% \author{...}
% \date{\today}
% \maketitle
% \tableofcontents
% \newpage

% --- Inhalt hier einfügen ---

\end{document}
  oder Inhalt direkt einfügen
%  Engine:      pdflatex (zwei Durchläufe für Inhaltsverzeichnis)
%               XeLaTeX nur bei Bedarf (fontspec/Spezialfonts)
% =====================================================================

\documentclass[11pt,a4paper,parskip=half]{scrartcl}

% --- Sprache & Eingabe ----------------------------------------------
\usepackage[utf8]{inputenc}
\usepackage[T1]{fontenc}
\usepackage[ngerman]{babel}
\usepackage{lmodern}
\usepackage{microtype}

% --- Mathematik -----------------------------------------------------
\usepackage{amsmath,amssymb,amsthm,mathtools}
\usepackage{cancel}                  % \cancel für Kürzen
\usepackage{siunitx}                 % \SI, \num, \si — deutsche Konvention
\sisetup{
  output-decimal-marker = {,},       % Komma statt Punkt
  per-mode              = symbol,    % m/s statt m·s^{-1}
  range-phrase          = {\,bis\,},
  list-final-separator  = { und },
  group-separator       = {\,},      % 10 000 mit schmalem Leerzeichen
}

% --- Layout ---------------------------------------------------------
\usepackage[a4paper,left=2.4cm,right=2.4cm,top=2.5cm,bottom=2.5cm]{geometry}
\usepackage{enumitem}
\setlist{topsep=2pt,itemsep=2pt}
\usepackage{array,booktabs,longtable,tabularx,multicol}
\usepackage{ragged2e}
\usepackage{fancyhdr}
\usepackage{hyperref}
\hypersetup{
  colorlinks = true,
  linkcolor  = mainblue,
  urlcolor   = mainblue,
  citecolor  = mainblue,
}

% --- Grafik ---------------------------------------------------------
\usepackage{tikz}
\usepackage{pgfplots}
\pgfplotsset{compat=1.18}
\usetikzlibrary{
  patterns, arrows.meta, calc, positioning,
  decorations.pathreplacing, decorations.markings,
  intersections, angles, quotes, 3d
}

% --- Chemie (nur einbinden, wenn benötigt) --------------------------
% \usepackage[version=4]{mhchem}     % Reaktionsgleichungen: \ce{H2O}
% \usepackage{chemfig}               % Strukturformeln
% \setatomsep{1.8em}
% \setbondoffset{1pt}
% \setdoublesep{0.45ex}

% =====================================================================
%  Farbschema (Kleeblatt-Konvention für Skripte)
% =====================================================================
\usepackage{xcolor}
\definecolor{mainblue}   {RGB}{ 0, 90,160}   % Definitionen, Hauptfarbe
\definecolor{darkblue}   {RGB}{ 0, 50,100}
\definecolor{exgreen}    {RGB}{ 34,120, 67}  % Beispiele
\definecolor{aufgorange} {RGB}{210,120, 20}  % Aufgaben
\definecolor{loesgray}   {RGB}{ 90, 90, 90}  % Lösungen
\definecolor{hinwyellow} {RGB}{200,160,  0}  % Hinweise
\definecolor{stolperred} {RGB}{180, 40, 40}  % Stolperfallen
\definecolor{tipppurple} {RGB}{120, 40,140}  % Klausurtipps
\definecolor{lighgray}   {RGB}{245,245,245}

% =====================================================================
%  Boxen-Umgebungen (tcolorbox)
% =====================================================================
\usepackage[most]{tcolorbox}
\tcbuselibrary{theorems, skins, breakable}

% --- Definition / Merksatz ------------------------------------------
\newtcolorbox{definition}[1][Definition]{%
  enhanced, breakable,
  colback=mainblue!7, colframe=mainblue,
  fonttitle=\bfseries,
  title=#1,
  attach boxed title to top left={yshift=-2mm,xshift=4mm},
  boxed title style={colback=mainblue,colframe=mainblue},
  arc=2pt, top=3mm, bottom=2mm
}

% --- Beispiel mit Musterlösung --------------------------------------
\newtcolorbox{beispiel}[1][Beispiel]{%
  enhanced, breakable,
  colback=exgreen!7, colframe=exgreen,
  fonttitle=\bfseries,
  title=#1,
  attach boxed title to top left={yshift=-2mm,xshift=4mm},
  boxed title style={colback=exgreen,colframe=exgreen},
  arc=2pt, top=3mm, bottom=2mm
}

% --- Übungsaufgabe --------------------------------------------------
\newtcolorbox{aufgabe}[1][Aufgabe]{%
  enhanced, breakable,
  colback=aufgorange!8, colframe=aufgorange,
  fonttitle=\bfseries,
  title=#1,
  attach boxed title to top left={yshift=-2mm,xshift=4mm},
  boxed title style={colback=aufgorange,colframe=aufgorange},
  arc=2pt, top=3mm, bottom=2mm
}

% --- Lösung ---------------------------------------------------------
\newtcolorbox{loesung}[1][Lösung]{%
  enhanced, breakable,
  colback=lighgray, colframe=loesgray,
  fonttitle=\bfseries,
  title=#1,
  attach boxed title to top left={yshift=-2mm,xshift=4mm},
  boxed title style={colback=loesgray,colframe=loesgray},
  arc=2pt, top=3mm, bottom=2mm
}

% --- Allgemeiner Hinweis --------------------------------------------
\newtcolorbox{hinweis}[1][Hinweis]{%
  enhanced, breakable,
  colback=hinwyellow!10, colframe=hinwyellow,
  fonttitle=\bfseries,
  title=#1,
  attach boxed title to top left={yshift=-2mm,xshift=4mm},
  boxed title style={colback=hinwyellow,colframe=hinwyellow},
  arc=2pt, top=3mm, bottom=2mm
}

% --- Stolperfalle (typischer Klausurfehler) -------------------------
\newtcolorbox{stolperfalle}[1][Stolperfalle]{%
  enhanced, breakable,
  colback=stolperred!8, colframe=stolperred,
  fonttitle=\bfseries\color{white},
  title=#1,
  attach boxed title to top left={yshift=-2mm,xshift=4mm},
  boxed title style={colback=stolperred,colframe=stolperred},
  arc=2pt, top=3mm, bottom=2mm
}

% --- Klausurtipp ----------------------------------------------------
\newtcolorbox{klausurtipp}[1][Klausurtipp]{%
  enhanced, breakable,
  colback=tipppurple!8, colframe=tipppurple,
  fonttitle=\bfseries\color{white},
  title=#1,
  attach boxed title to top left={yshift=-2mm,xshift=4mm},
  boxed title style={colback=tipppurple,colframe=tipppurple},
  arc=2pt, top=3mm, bottom=2mm
}

% =====================================================================
%  Sonstige Konventionen
% =====================================================================
\renewcommand{\vec}[1]{\boldsymbol{#1}}        % Vektoren fett statt Pfeil
% Alternative: \renewcommand{\vec}[1]{\overrightarrow{#1}}

\newcommand{\R}{\mathbb{R}}
\newcommand{\N}{\mathbb{N}}
\newcommand{\Z}{\mathbb{Z}}
\newcommand{\Q}{\mathbb{Q}}

% Endergebnis einrahmen (für Aufgabenlösungen)
\newcommand{\ergebnis}[1]{\boxed{\;#1\;}}

% Kurzbefehl für Bildungsplan-Code-Verweise
\newcommand{\bpref}[1]{\textit{(Bildungsplan #1)}}

% =====================================================================
%  Kopf- und Fußzeile (optional, hier aktiv)
% =====================================================================
\pagestyle{fancy}
\fancyhf{}
\fancyhead[L]{\small\itshape\leftmark}
\fancyhead[R]{\small Klasse 10 -- Gymnasium BW}
\fancyfoot[C]{\thepage}
\renewcommand{\headrulewidth}{0.4pt}

% =====================================================================
%  Beginn des Dokuments
% =====================================================================
\begin{document}

% \title{...}
% \author{...}
% \date{\today}
% \maketitle
% \tableofcontents
% \newpage

% --- Inhalt hier einfügen ---

\end{document}
